% GENERATED FILE. Edit canonical lessons, then run npm run generate:books.
% canonical-source-hash: fnv1a64:14636abe69a09bb6
% canonical-lessons: SA-C210-acaryah, SA-C210-pradhanacaryah, SA-C210-vidyarthi, SA-C210-sahapathi, SA-C210-cikitsakah

\chapter{Preceptor, Head teacher, Pupil, Classmate, Doctor}
\label{ch:sa-preceptor-head-teacher-pupil-classmate-doctor}
\hlchaptermodality{210}

\begin{chapteropening}
\textbf{By the end of this chapter:} \emph{I can say a preceptor, a head teacher, a pupil, a classmate and a doctor.}

Complete the last lesson of chapter 210: I can say a preceptor, a head teacher, a pupil, a classmate and a doctor.
\end{chapteropening}

\section[\texorpdfstring{\sk{आचार्यः} (ācāryaḥ)}{आचार्यः (ācāryaḥ)}]{\sk{आचार्यः} (ācāryaḥ) — a preceptor, a teacher}

\label{lesson:SA-C210-acaryah}



\begin{quote}
Before the new one: say the Sanskrit for a painter, then the Sanskrit for a craftsman.
\end{quote}

\subsection*{You'll want to know: \sk{आचार्यः}}
\textbf{\sk{आचार्यः}} — \emph{ācāryaḥ} — \textquotedblleft{}a preceptor, a teacher\textquotedblright{}.

\begin{grammarlens}[title={What you've built}]
One more word for this chapter.
\end{grammarlens}

\subsection*{Your turn}
\begin{itemize}
\raggedright
  \item \emph{Say it:} \emph{ācāryaḥ}
  \item \emph{Say it:} \emph{ācāryaḥ}, once more
  \item \emph{Say it:} say \emph{ācāryaḥ}
  \item \emph{Recall:} say the Sanskrit for a painter, then the Sanskrit for a craftsman, then say \emph{ācāryaḥ} again
\end{itemize}

\begin{tcolorbox}[breakable,colback=teal!4,colframe=teal!35!black,title={Before you move on}]
What does \sk{आचार्यः} mean? (\textquotedblleft{}A preceptor, a teacher\textquotedblright{}.) Say it once more.
\end{tcolorbox}
\section[\texorpdfstring{\sk{प्रधानाचार्यः} (pradhānācāryaḥ)}{प्रधानाचार्यः (pradhānācāryaḥ)}]{\sk{प्रधानाचार्यः} (pradhānācāryaḥ) — a head teacher}

\label{lesson:SA-C210-pradhanacaryah}



\begin{quote}
Before the new one: say the Sanskrit for a craftsman, then the Sanskrit for a preceptor.
\end{quote}

\subsection*{You'll want to know: \sk{प्रधानाचार्यः}}
\textbf{\sk{प्रधानाचार्यः}} — \emph{pradhānācāryaḥ} — \textquotedblleft{}a head teacher\textquotedblright{}.

\begin{grammarlens}[title={What you've built}]
One more word for this chapter.
\end{grammarlens}

\subsection*{Your turn}
\begin{itemize}
\raggedright
  \item \emph{Say it:} \emph{pradhānācāryaḥ}
  \item \emph{Say it:} \emph{pradhānācāryaḥ}, once more
  \item \emph{Say it:} say \emph{pradhānācāryaḥ}
  \item \emph{Recall:} say the Sanskrit for a craftsman, then the Sanskrit for a preceptor, then say \emph{pradhānācāryaḥ} again
\end{itemize}

\begin{tcolorbox}[breakable,colback=teal!4,colframe=teal!35!black,title={Before you move on}]
What does \sk{प्रधानाचार्यः} mean? (\textquotedblleft{}A head teacher\textquotedblright{}.) Say it once more.
\end{tcolorbox}
\section[\texorpdfstring{\sk{विद्यार्थी} (vidyārthī)}{विद्यार्थी (vidyārthī)}]{\sk{विद्यार्थी} (vidyārthī) — a pupil, a learner}

\label{lesson:SA-C210-vidyarthi}



\begin{quote}
Before the new one: say the Sanskrit for a preceptor, then the Sanskrit for a head teacher.
\end{quote}

\subsection*{You'll want to know: \sk{विद्यार्थी}}
\textbf{\sk{विद्यार्थी}} — \emph{vidyārthī} — \textquotedblleft{}a pupil, a learner\textquotedblright{}.

\begin{grammarlens}[title={What you've built}]
One more word for this chapter.
\end{grammarlens}

\subsection*{Your turn}
\begin{itemize}
\raggedright
  \item \emph{Say it:} \emph{vidyārthī}
  \item \emph{Say it:} \emph{vidyārthī}, once more
  \item \emph{Say it:} say \emph{vidyārthī}
  \item \emph{Recall:} say the Sanskrit for a preceptor, then the Sanskrit for a head teacher, then say \emph{vidyārthī} again
\end{itemize}

\begin{tcolorbox}[breakable,colback=teal!4,colframe=teal!35!black,title={Before you move on}]
What does \sk{विद्यार्थी} mean? (\textquotedblleft{}A pupil, a learner\textquotedblright{}.) Say it once more.
\end{tcolorbox}
\section[\texorpdfstring{\sk{सहपाठी} (sahapāṭhī)}{सहपाठी (sahapāṭhī)}]{\sk{सहपाठी} (sahapāṭhī) — a classmate}

\label{lesson:SA-C210-sahapathi}



\begin{quote}
Before the new one: say the Sanskrit for a head teacher, then the Sanskrit for a pupil.
\end{quote}

\subsection*{You'll want to know: \sk{सहपाठी}}
\textbf{\sk{सहपाठी}} — \emph{sahapāṭhī} — \textquotedblleft{}a classmate\textquotedblright{}.

\begin{grammarlens}[title={What you've built}]
One more word for this chapter.
\end{grammarlens}

\subsection*{Your turn}
\begin{itemize}
\raggedright
  \item \emph{Say it:} \emph{sahapāṭhī}
  \item \emph{Say it:} \emph{sahapāṭhī}, once more
  \item \emph{Say it:} say \emph{sahapāṭhī}
  \item \emph{Recall:} say the Sanskrit for a head teacher, then the Sanskrit for a pupil, then say \emph{sahapāṭhī} again
\end{itemize}

\begin{tcolorbox}[breakable,colback=teal!4,colframe=teal!35!black,title={Before you move on}]
What does \sk{सहपाठी} mean? (\textquotedblleft{}A classmate\textquotedblright{}.) Say it once more.
\end{tcolorbox}
\section[\texorpdfstring{\sk{चिकित्सकः} (cikitsakaḥ)}{चिकित्सकः (cikitsakaḥ)}]{\sk{चिकित्सकः} (cikitsakaḥ) — a doctor}

\label{lesson:SA-C210-cikitsakah}



\begin{quote}
Before the new one: say the Sanskrit for a pupil, then the Sanskrit for a classmate.
\end{quote}

\subsection*{You'll want to know: \sk{चिकित्सकः}}
\textbf{\sk{चिकित्सकः}} — \emph{cikitsakaḥ} — \textquotedblleft{}a doctor\textquotedblright{}.

\begin{grammarlens}[title={What you've built}]
One more word for this chapter.
\end{grammarlens}

\subsection*{Your turn}
\begin{itemize}
\raggedright
  \item \emph{Say it:} \emph{cikitsakaḥ}
  \item \emph{Say it:} \emph{cikitsakaḥ}, once more
  \item \emph{Say it:} say \emph{cikitsakaḥ}
  \item \emph{Recall:} say the Sanskrit for a pupil, then the Sanskrit for a classmate, then say \emph{cikitsakaḥ} again
\end{itemize}

\begin{tcolorbox}[breakable,colback=teal!4,colframe=teal!35!black,title={Before you move on}]
What does \sk{चिकित्सकः} mean? (\textquotedblleft{}A doctor\textquotedblright{}.) Say it once more.
\end{tcolorbox}
